\section{Cyclic markings and exact presentations}\label{sec:cyclic}
\subsection{The precise input}
For a group generated by a finite named set $S$, write
$A(R)=\langle r:r\in R\rangle$ for $R\subseteq S$, including
$A(\varnothing)=1$. We call this a parabolic marking if
\begin{equation}\label{C-eq:ip}
 A(R)\cap A(Q)=A(R\cap Q)
\end{equation}
and the cosets of all $A(R)$ have the finite Helly property: every finite
pairwise-intersecting collection has a common element. The names denote
distinct nonidentity elements. The standard coset system has maximal
parabolics $A(S\setminus\{s\})$.

Here are the classical inputs retained for an arbitrary Artin--Tits
marking. Standard parabolic presentations embed, their intersections obey
(\ref{C-eq:ip}), and their cosets have this Helly property. These are the
standard-parabolic results invoked in N005 Theorems~2.11--2.12:
van der Lek's parabolic embedding and the intersection/product results
cited there from Godelle--Paris. In particular, each singleton Artin
parabolic is infinite cyclic. For a Coxeter marking the corresponding
facts are the classical Coxeter parabolic theorem, and singleton
parabolics have order two. We do not reprove those external
results or identify their proof costs with zero.
Precisely, the source locators are Bourbaki, \emph{Groupes et alg\`ebres
de Lie}, Chapters 4--6 (2006 edition), Chapter 4, \S1, Theorem 2
and Exercise \S1.1(b); H.~van der Lek, \emph{The homotopy type of
complex hyperplane complements} (1983), Chapter II, Theorem 4.13;
and E.~Godelle--L.~Paris, \emph{$K(\pi,1)$ and word problems for
infinite type Artin--Tits groups, and applications to virtual braid
groups}, Math.~Z.~272 (2012), Proposition 4.5. We retain these
interfaces as the primary source does, rather than claiming fresh
proof verification of the entire cited works.

Lemma~\ref{F-lem:cosets} identifies this Helly hypothesis with flag transitivity, once formula (\ref{C-eq:ip}) identifies generated supports with standard intersections.

\begin{lemma}\label{C-lem:product}
In a parabolic marking, for $H=A(R)$, $K=A(Q)$ and $L=A(P)$,
\[
 HK\cap HL=H(K\cap L).
\]
The complex of cosets of maximal parabolics, with simplices the families
having common intersection, is a flag complex. Its coset geometry is
residually connected; its panel of type $s$ has $|\langle s\rangle|$
chambers.
\end{lemma}
\begin{proof}
If $z\in HK\cap HL$, the left cosets $H,zK,zL$ meet pairwise:
the first meets each of the others and the latter two contain $z$.
Helly gives $h=zk=zl\in H$, with $k\in K,l\in L$; hence
$k=l\in K\cap L$ and $z\in H(K\cap L)$. The reverse inclusion
is immediate. The complex assertion is exactly Helly. Every standard
residue has acting group $A(R)$ and maximal parabolics
$A(R\setminus\{s\})$, $s\in R$. If $|R|\geq2$ they generate
$A(R)$: each named generator belongs to a child missing a different
name. Its incidence graph is therefore connected (successive
multiplications by parabolic elements give incidence paths). Flag
transitivity identifies every residue with a translate of one of these.
The Borel is $1$ by (\ref{C-eq:ip}), and the rank-one residue is
$\langle s\rangle/1$, proving the panel assertion.
\end{proof}

\subsection{One cyclic marking theorem}
Let $S=Y\sqcup L$ with
$L=\{x_0,\ldots,x_{k-1}\}$ and $k\geq2$. Suppose $(A,S)$ is a
parabolic marking and $\sigma\in\Aut(A)$ fixes every $y\in Y$ and
sends $x_j$ to $x_{j+1}$, with subscripts modulo $k$.
Then $\sigma^k=\id$ since it fixes every generator. Set
\[
 G=A\rtimes C_k=\langle A,t\mid t^k=1,\ t^{-1}at=\sigma(a)
 \ (a\in A)\rangle,
 \qquad T=Y\sqcup\{x_0,t\}.
\]
This notation fixes the conjugation convention. The quotient onto $C_k$
proves that $t$ has exact order $k$. The set $T$ generates $G$, since
$x_j=t^{-j}x_0t^j$.

For $R\subseteq T$ define a subset $E(R)$ of $S$ by
\[
 E(R)=\begin{cases}
 R,&t\notin R,\\
 R\cap Y,&t\in R,\ x_0\notin R,\\
 (R\cap Y)\cup L,&t,x_0\in R.
 \end{cases}
\]
In the first line $R$ is read inside $Y\sqcup\{x_0\}\subseteq S$.
\begin{theorem}\label{C-thm:cyclic}
For every $R\subseteq T$ the subgroup generated by $R$ is
\begin{equation}\label{C-eq:all-subgroups}
 G(R)=\begin{cases}
 A(E(R)),&t\notin R,\\
 A(E(R))\rtimes C_k,&t\in R.
 \end{cases}
\end{equation}
The marking $(G,T)$ has the intersection and coset Helly properties.
Consequently its standard coset system is flag-transitive, residually
connected and firm, and its coset complex is flag. The product identity
of Lemma~\ref{C-lem:product} holds for every three generated parabolics.
Its panels have size $k$ at $t$ and $|\langle s\rangle|$ at every
$s\in Y\sqcup\{x_0\}$.
\end{theorem}
\begin{proof}
If $t\notin R$, the assertion is the embedding of $A$ in $G$.
If $t\in R$ but $x_0\notin R$, it fixes every named generator in
$R\cap Y$, giving the second case. If both belong to $R$, conjugating
$x_0$ produces exactly the whole orbit $L$, and this subgroup of $A$
is invariant under $t$. This proves (\ref{C-eq:all-subgroups}) without
an order calculation or finiteness assumption on $A$.

If both subsets contain $t$, the two $A$-parts intersect by
(\ref{C-eq:ip}); the orbit $L$ is present in their intersection precisely
when $x_0$ belongs to both subsets. Thus their intersection is the
subgroup for $R\cap Q$. If neither contains $t$, use (\ref{C-eq:ip})
inside $A$. In the mixed case an element of the intersection has trivial
$C_k$ coordinate. The smaller support contains at most $x_0$ from $L$,
so intersecting its support with the other support again gives exactly
$E(R\cap Q)$. This proves full generated-subgroup intersection.

It remains to prove Helly, which does not follow from subgroup
intersection alone. The maximal parabolics just computed are
\begin{equation}\label{C-eq:maximal}
 P_t=A(Y\cup\{x_0\}),\quad
 P_{x_0}=A(Y)\rtimes C_k,\quad
 P_y=A(S\setminus\{y\})\rtimes C_k\ (y\in Y).
\end{equation}
The $A$-supports in the latter groups are invariant under $C_k$.
Take pairwise-intersecting left cosets of these groups. Repeated types
coincide. If a $P_t$ coset occurs, translate its $C_k$ coordinate to
$1$. All other cosets project onto $C_k$ and have representatives
$(a,1)$. Their fibers at coordinate $1$ are cosets of the indicated
standard parabolics in $A$. Two cosets of the latter kind intersect
if and only if their $A$ fibers intersect, because they are Cartesian
products with $C_k$; incidence with the $P_t$ coset likewise takes
place in this fiber. Thus the $A$ fibers meet pairwise and base Helly
gives a common representative. If no $P_t$ coset occurs, the same
argument applies at coordinate $1$ without translation. This proves
Helly for maximal parabolics. Replace a coset of any generated
parabolic by its maximal-parabolic cosets, using the intersection
property just proved, to get Helly for all of them. Finally apply
Lemma~\ref{C-lem:product}. Every named generator is nontrivial and the
panel sizes are at least two, proving firmness.
\end{proof}
This is the rank-one cyclic case of the orbit construction used in
Piedade--Tranchida's papers on twisting. Its proof is included in full
so that its geometric and presentation conclusions have a complete local proof. It does not assert that an arbitrary Shephard marking has
these properties.

\subsection{The involutory presentation and the doubled braid}
Write $\W_m(u,v)$ for the alternating word of \emph{length} $m$,
starting with $u$. It is not the ordinary power $(uv)^m$.
Let $A$ be the Artin group on $Y\sqcup\{x,x'\}$, with Artin
relations $\mathcal R_Y$ on $Y$, identical labels $m_y$ on
$(x,y)$ and $(x',y)$, and label $m$ on $(x,x')$.
All labels lie in $\{2,3,\ldots\}\cup\{\infty\}$; an infinite
label means no relation. Interchanging $x,x'$ is an automorphism of
the presented group, fixes $Y$, and has order two because the names
are distinct by the retained Artin parabolic theorem.

\begin{lemma}\label{C-lem:braid}
If $t^2=1$, then for every $m\geq1$,
\[
 \W_m(x,txt)=(xt)^m t^{m\bmod2},\qquad
 \W_m(txt,x)=(tx)^m t^{m\bmod2}.
\]
Consequently the braid of length $m$ on $x,txt$ is equivalent to the
braid of length $2m$ on $x,t$.
\end{lemma}
\begin{proof}
For $m=2r$, the two words are $(xtxt)^r=(xt)^{2r}$ and
$(txtx)^r=(tx)^{2r}$. For $m=2r+1$, append respectively $x$ and
$txt$ to those words, obtaining $(xt)^{2r+1}t$ and
$(tx)^{2r+1}t$. Cancellation of the common final $t$ in the odd case
proves the equivalence. The alternating words of length $2m$ are
exactly $(xt)^m$ and $(tx)^m$.
\end{proof}

\begin{theorem}\label{C-thm:fold}
The semidirect product which interchanges $x,x'$ has presentation
\begin{align*}
 \langle Y,x,t\mid {}&\mathcal R_Y,\ t^2=1,\ [t,y]=1\ (y\in Y),\\
 &\W_{m_y}(x,y)=\W_{m_y}(y,x)\ (m_y<\infty),\\
 &\W_{2m}(x,t)=\W_{2m}(t,x)\ (m<\infty)\rangle.
\end{align*}
This is a Shephard presentation, and its named marking satisfies every
conclusion of Theorem~\ref{C-thm:cyclic}.
\end{theorem}
\begin{proof}
Map the displayed generators to the semidirect product. They satisfy
every relator: the last by Lemma~\ref{C-lem:braid}; the others follow
from the defining action and the original Artin relations. They generate,
since $x'=txt$. Conversely, in the displayed presentation set
$x'=txt$. Conjugating a mixed braid by $t$, which fixes $y$, gives
the braid on $x',y$. Lemma~\ref{C-lem:braid} supplies the braid on
$x,x'$. The relations on $Y$ are unchanged. Thus the original Artin
presentation maps into this group, and conjugation by $t$ implements
its transposition of the two leaves. The semidirect-product universal
property gives the inverse map. Both composites fix all named
 generators, so these maps are inverse. This establishes the exact
 group and marking, not merely a quotient action. The cyclic theorem
with $k=2$ now applies. An edge absent from the drawing has label two,
so all commutations $[t,y]=1$ remain part of the presentation.
\end{proof}

\subsection{Cyclic amalgamation}
Let $A_0$ be an Artin group on $Y\sqcup\{x\}$ and
$C=A_0(Y)$. The parabolic embedding theorem identifies $C$ with the
Artin group on the induced graph on $Y$. Form the amalgam $M$ of
$k\geq2$ copies of $A_0$ along this common $C$. Its presentation has
$Y$ and $x_0,\ldots,x_{k-1}$, the relations $\mathcal R_Y$, and the
same braid between $x_j$ and each $y\in Y$. There are no relations
between distinct $x_j$. This follows directly from the universal
property of the amalgam; the embeddings of each factor are the
classical normal-form theorem. The resulting presentation is itself
Artin, with all distinct leaf-pair labels infinite.

\begin{theorem}\label{C-thm:amalgam}
Cyclic permutation of the copies defines an automorphism $\sigma$ of
$M$. With $t^{-1}x_jt=x_{j+1}$, its semidirect product has presentation
\[
 M\rtimes C_k\cong
 \langle Y,x_0,t\mid\mathcal R_Y,\ t^k=1,\ [t,y]=1\ (y\in Y),\
 \W_{m_y}(x_0,y)=\W_{m_y}(y,x_0)\ (m_y<\infty)\rangle.
\]
There is no relation on the pair $x_0,t$. This Shephard marking has
full generated-parabolic intersection and the product identity,
and its standard coset system is flag-transitive, residually connected
and firm with flag coset complex. Its panels have sizes $k$ and
infinity.
\end{theorem}
\begin{proof}
The permutation preserves the Artin relators and has order $k$ on
the distinct leaf names. The displayed generators map to the
semidirect product and generate it. In the converse direction define
$x_j=t^{-j}x_0t^j$. Since $t$ fixes $Y$, conjugation by $t^j$
transports every mixed braid to the required braid on $x_j,y$.
The relation $t^k=1$ closes the cycle. The Artin presentation of
$M$ and then its semidirect product therefore map into the displayed
group. Both composites fix $Y,x_0,t$, proving inverse maps. Apply
Theorem~\ref{C-thm:cyclic} to the Artin marking of $M$.
\end{proof}
